\subsection{拓扑空间的逆极限}

设 I 是一个偏序（但不一定有向）集（*），其上的序关系记作 \(\alpha \leqslant \beta\)。对每个 \(\alpha \in I\)，设 \(X_{\alpha}\) 是一个拓扑空间，且对每一对 \((\alpha, \beta)\)（满足 \(\alpha \leqslant \beta\)），设 \(f_{\alpha\beta}\) 是 \(X_{\beta}\) 到 \(X_{\alpha}\) 中的一个映射。我们说 \((X_{\alpha}, f_{\alpha\beta})\) 是拓扑空间的一个逆系统，如果：1) \((X_{\alpha}, f_{\alpha\beta})\) 是集合的一个逆系统；2) \(f_{\alpha\beta}\) 是连续映射（每当 \(\alpha \leqslant \beta\)）。设 X 表示集合 \(\varprojlim X_{\alpha}\)，且对每个 \(\alpha \in I\)，设 \(f_{\alpha}\) 是典范映射 \(X \to X_{\alpha}\)；则 X 上使诸 \(f_{\alpha}\) 连续的最粗的拓扑称为（关于 \(f_{\alpha\beta}\) 的）诸 \(X_{\alpha}\) 的拓扑的逆极限，而赋有该拓扑的集合 X 称为拓扑空间的逆系统 \((X_{\alpha}, f_{\alpha\beta})\) 的逆极限。每当我们把 \(\varprojlim X_{\alpha}\) 作为拓扑空间来谈论时，总应理解成这个空间的拓扑是诸 \(X_{\alpha}\) 的拓扑的逆极限，除非另有明确说明。

集合 X 是积 \(\prod_{\alpha\in I}X_{\alpha}\) 的由满足下列条件的点 x 组成的子集：

\[
\operatorname{pr} _ {\alpha} (x) = f _ {\alpha \beta} (\operatorname{pr} _ {\beta} (x))\tag{I}
\]

每当 \(\alpha \leqslant \beta\)。由第 1 目命题 3 可知，诸 \(X_{\alpha}\) 的拓扑的逆极限与 \(X\) 上由积空间 \(\prod_{\alpha \in I} X_{\alpha}\) 的拓扑诱导的拓扑相同。若对每个 \(\alpha \in I\)，\(Y_{\alpha}\) 是 \(X_{\alpha}\) 的子空间，且诸 \(Y_{\alpha}\) 构成诸 \(X_{\alpha}\) 的子集的一个逆系统（《集合论》，第三章，§ 7，第 2 目），则显然拓扑空间 \(\varprojlim Y_{\alpha}\) 是 \(\varprojlim X_{\alpha}\) 的一个子空间。

设 \((\mathbf{X}_{\alpha}^{\prime}, f_{\alpha\beta}^{\prime})\) 是以同一集合 I 为指标集的拓扑空间的另一个逆系统，且对每个 \(\alpha \in I\)，设 \(u_{\alpha}: X_{\alpha} \to X_{\alpha}^{\prime}\) 是连续映射，使得 \((u_{\alpha})\) 是映射的一个逆系统；则 \(u = \varprojlim u_{\alpha}\) 是 \(X = \varprojlim X_{\alpha}\) 到 \(X^{\prime} = \varprojlim X_{\alpha}^{\prime}\) 中的一个连续映射。因为若 \(f_{\alpha}^{\prime}\) 是典范映射 \(X^{\prime} \to X_{\alpha}^{\prime}\)，我们有 \(f_{\alpha}^{\prime} \circ u = u_{\alpha} \circ f_{\alpha}\)，从而 \(f_{\alpha}^{\prime} \circ u\) 对每个 \(\alpha \in I\) 连续，于是断言由 § 2 第 3 目命题 4 得出。

最后，设 I 是有向集，并设 J 是 I 的一个共尾子集；设 Z 是拓扑空间的逆系统 \((\mathbf{X}_{\alpha}, f_{\alpha\beta})_{\alpha\in\mathbf{J}, \beta\in\mathbf{J}}\) 的逆极限。则典范双射 \(g: X \to Z\)（《集合论》，第三章，§ 7，第 2 目，命题 3）是一个同胚。因为我们有 \(\operatorname{pr}_{\lambda}(g(x)) = \operatorname{pr}_{\lambda}(x)\)（对每个 \(\lambda \in J\)）；故 g 连续（第 1 目，命题 1）；又若 h 是 g 的逆，则对每个 \(\alpha \in I\)，存在 \(\lambda \in J\) 使 \(\alpha \leqslant \lambda\)，从而 \(\operatorname{pr}_{\alpha}(h(z)) = f_{\alpha\lambda}(\operatorname{pr}_{\lambda}(z))\)；由于诸 \(f_{\alpha\lambda}\) 连续，这表明 h 连续（第 1 目，命题 1）。

命题 9. 设 I 是有向集，J 是 I 的一个共尾子集。设 \((\mathbf{X}_{\alpha}, f_{\alpha \beta})\) 是以 I 为指标集的拓扑空间的一个逆系统；设 \(\mathbf{X} = \varprojlim \mathbf{X}_{\alpha}\)，并设 \(f_{\alpha}: \mathbf{X} \to \mathbf{X}_{\alpha}\) 是典范映射。则由诸集合 \(\overline{f}_{\alpha}^{1}(\mathbf{U}_{\alpha})\) 组成的族（其中 \(\alpha\) 取遍 J，\(\mathbf{U}_{\alpha}\) 取遍基 \(\mathfrak{B}_{\alpha}\)——\(\mathbf{X}_{\alpha}\) 的拓扑的基，对每个 \(\alpha \in \mathbf{J}\)）是 \(\mathbf{X}\) 的拓扑的一个基。

由 § 2 第 3 目可知，形如 \(\overline{f}_{\alpha}^{1}(\mathbf{U}_{\alpha})\)（\(\alpha \in \mathbf{I}\)，\(\mathbf{U}_{\alpha}\) 是 \(\mathbf{X}_{\alpha}\) 中的开集）的集合的有限交组成 \(\mathbf{X}\) 的拓扑的一个基。若 \((\alpha_i)_{1 \leqslant i \leqslant n}\) 是 \(\mathbf{I}\) 的一个有限指标族，则存在 \(\gamma \in \mathbf{J}\)，使 \(\alpha_i \leqslant \gamma\)（对 \(1 \leqslant i \leqslant n\)）；从而 \(f_{\alpha_i} = f_{\alpha_i \gamma} \circ f_\gamma\)；若令

\[
\begin{array}{c} \mathrm{V} _ {\gamma} = \bigcap_ {i} \overline {{f}} _ {\alpha_ {i} \gamma} ^ {1} (\mathrm{U} _ {\alpha_ {i}}), \\ \overline {{f}} _ {\gamma} ^ {1} (\mathrm{V} _ {\gamma}) = \bigcap_ {i} \overline {{f}} _ {\alpha_ {i}} ^ {1} (\mathrm{U} _ {\alpha_ {i}}); \end{array}
\]

则我们有

而 \(V_{\gamma}\) 是开集，因而是属于 \(\mathfrak{B}_{\gamma}\) 的一些集合之并。由此即得结论。

推论. 设 A 是 X 的一个子集，且用 \(A_{\alpha}\) 表示 \(f_{\alpha}(A)\)（对每个 \(\alpha \in I\)）。则：

\begin{enumerate}
\def\labelenumi{(\roman{enumi})}
\item
  诸 \(\mathbf{A}_{\alpha}\)（相应地，诸 \(\overline{\mathbf{A}}_{\alpha}\)）构成诸 \(\mathbf{X}_{\alpha}\) 的子集的一个逆系统，且 \(\overline{\mathbf{A}} = \bigcap_{\alpha} \overline{f}_{\alpha}^{1}(\overline{\mathbf{A}}_{\alpha}) = \varprojlim \overline{\mathbf{A}}_{\alpha}\)。
\item
  若 A 在 X 中是闭集，则 \(\mathbf{A} = \varprojlim \mathbf{A}_{\alpha} = \varprojlim \overline{\mathbf{A}}_{\alpha}\)。
\end{enumerate}

(i) 的第一个断言由诸关系 \(f_{\alpha} = f_{\alpha \beta} \circ f_{\beta}\)（对 \(\alpha \leqslant \beta\)）及诸 \(f_{\alpha \beta}\) 的连续性得出（§ 2，第 1 目，定理 1）。用 \(A'\) 表示

\[
\bigcap_ {\alpha} \overline {{f}} _ {\alpha} ^ {1} (\overline {{\mathrm{A}}} _ {\alpha});
\]

则显然 \(\mathbf{A}'\) 是闭集且包含 \(\mathbf{A}\)，从而 \(\overline{\mathbf{A}} \subset \mathbf{A}'\)。反之，设 \(x \in \mathbf{A}'\)；我们必须证明 \(x\) 属于 \(\mathbf{A}\) 的闭包。凭借命题 9，只须证明 \(x\) 的每个形如 \(\overline{f}_{\alpha}^{1}(\mathrm{U}_{\alpha})\) 的邻域（其中 \(\alpha \in \mathbf{I}\)，\(\mathbf{U}_{\alpha}\) 是 \(\mathbf{X}_{\alpha}\) 中的开集）都与 \(\mathbf{A}\) 相交。而按假设，\(f_{\alpha}(x) \in \mathbf{U}_{\alpha}\)，又因 \(f_{\alpha}(x) \in \overline{\mathbf{A}}_{\alpha}\)，我们有 \(\mathbf{U}_{\alpha} \cap \mathbf{A}_{\alpha} \neq \emptyset\)，这就意味着 \(\mathbf{A} \cap \overline{f}_{\alpha}^{1}(\mathbf{U}_{\alpha})\) 非空。

为建立 (ii)，只须注意，对 A 不加任何限制时，我们有 \(A \subset \varprojlim A_{\alpha} \subset \varprojlim \overline{A}_{\alpha}\)；而若 A 是闭集，则由 (i)

\[
\mathrm{A} = \varprojlim \overline {{\mathrm{A}}} _ {\alpha}
\]

便得 (ii)。

例. 设 I 是有向集，\((\mathbf{X}_{\alpha})_{\alpha \in I}\) 是集合 Y 的子集的一个族，使得 \(\mathbf{X}_{\alpha} \supset \mathbf{X}_{\beta}\)（每当 \(\alpha \leqslant \beta\)）。对每个 \(\alpha \in I\)，设 \(\mathfrak{G}_{\alpha}\) 是 \(\mathbf{X}_{\alpha}\) 上的一个拓扑，使得 \(\mathfrak{G}_{\beta}\) 都细于 \(\mathbf{X}_{\beta}\) 上由 \(\mathfrak{G}_{\alpha}\) 诱导的拓扑（每当 \(\alpha \leqslant \beta\)）。若取 \(f_{\alpha \beta}\) 为典范单射 \(\mathbf{X}_{\beta} \to \mathbf{X}_{\alpha}\)（对 \(\alpha \leqslant \beta\)），则 \(\varprojlim \mathbf{X}_{\alpha}\) 可以典范地等同于诸 \(\mathbf{X}_{\alpha}\) 的交 X，赋以诸拓扑之上确界（§ 2，第 3 目，例 2）——这些拓扑由诸 \(\mathfrak{G}_{\alpha}\) 在 X 上诱导。

